\section{Prácticas típicas de GenAI empresarial}

\subsection{Asistentes inteligentes en los servicios financieros}

La ruta habitual de GenAI en la banca digital consiste en integrar el LLM en el servicio al cliente, la gestión patrimonial y los procesos internos de cumplimiento normativo. El patrón típico es el siguiente:

\begin{itemize}
    \item \textbf{Entrada}: el usuario pregunta por asesoría de inversión / reporta una transacción sospechosa
    \item \textbf{Búsqueda}: recopila en tiempo real datos de mercado y normativa regulatoria
    \item \textbf{LLM}: genera una respuesta explicativa más una recomendación de acción
    \item \textbf{Herramientas}: llena automáticamente los reportes y genera borradores de contratos
\end{itemize}

\begin{knowledgebox}{El ciclo cerrado de tres capas de la IA financiera}
1. Frontend: el LLM genera una respuesta explicativa<br>
2. Capa intermedia: la búsqueda/recuperación dinámica garantiza que la información esté actualizada<br>
3. Backend: herramientas de automatización de bajo código ejecutan las instrucciones<br>
Entre las tres capas debe garantizarse la trazabilidad (¿por qué se recomienda esta estrategia?) y la prevención de riesgos (rechazar operaciones sensibles).
\end{knowledgebox}

\subsection{El GenAI Agent en la manufactura y la cadena de suministro}

Las empresas manufactureras usan el GenAI Agent para conectar el MES, el ERP y las herramientas de visualización de la cadena de suministro:

\begin{itemize}
    \item El ingeniero sube el CAD junto con la descripción del caso de uso
    \item El Agent lee el BOM, los registros de prueba y el historial de fallas, y genera un plan de diagnóstico de fallas
    \item Activa directamente una orden de trabajo o notifica la planeación de compras
\end{itemize}

\begin{warningbox}{El escenario de prueba debe reproducir la realidad}
Muchos agentes se desempeñan bien en el entorno de prueba, pero fallan en cuanto entran a un taller real debido a la diferencia en la distribución de los datos (bitácoras frente a sensores en tiempo real). Es necesario recopilar feedback durante la operación real y reproducirlo para el modelo.
\end{warningbox}

\subsection{Resumen de la sección}

Las prácticas de los distintos sectores siguen el mismo patrón de ciclo cerrado: la IA aporta comprensión y generación, la búsqueda y los datos garantizan la exactitud factual, y las herramientas y procesos completan la ejecución; toda la cadena requiere trazabilidad y control.

